\chapter{19}


Während Louis die Wiese durchquerte, sah er sie vor sich. Seine
«Taylor». Die erste selbst finanzierte, richtig gute Akustikgitarre,
zusammengespart über viele Monate. Wo mochte sie sein?
Sie waren eins geworden über die Zeit, das Instrument und er.
An eine Neue würde er sich gewöhnen müssen. Zu grosse Wünsche waren
deplatziert. Sein Budget begrenzt. 
Als Erstes würde er \emph{«L’amore è»} spielen, sich Giorgia widmen, den innigen Stunden mit ihr und ihrem gemeinsamen Song. Dem einzigen. Die Wehmut war zurück. Er würde sich die Momente in Erinnerung rufen, diese ersten, leichten. Sich an Giorgias Stimme erinnern, daran, wie sie zusammen den Text gelernt hatten.

\qrblock{Enrico Nigiotti \& Matilde Gioli \emph{«L’amore è»}}{media/QR-Codes/019_L`amore_è_Enrico_Nigiotti_&_Matilde_Gioli_Youtube.png}

Warum er diesen Schmerz nur immer wieder befeuerte? Es war wohl
einfach zu schön gewesen. Damals, als sie nah beisammen gesessen und
sich herausgefordert hatten. Liebevoll. Spassig auch. 
Giorgia hatte sich erst geziert. Und sich dann mitten im Singen vor
Lachen geschüttelt. Louis hatte die Mundwinkel hinuntergezogen und den
Beleidigten inszeniert. Giorgia ergab sich, was selten vorkam, und
«L’amore è» war zu einem Song geworden, der schliesslich nur
ihnen beiden gehörte.
«Credimi, credimi, credimi, credimi, credimi sempre.»
Zuweilen hatten sie die Worte vor Freude rausgebrüllt.
Je mehr sie getrunken hatten, desto leichter verhaspelten sie sich und
\emph{«credimi»} wurde zu hastigem Wortschnappen. Beide genossen sie
ihre Heiterkeit. Zusammen waren sie ein fulminantes Feuerwerk.
Es hatte keine Zweifel gegeben, auf beiden Seiten. Natürlich glaubten
sie einander alles und immer, dachten sie und waren sich sicher. Zu
dieser Zeit.

Der Song blieb der Einzige, den Giorgia mitsang. Und auch nur, wenn sie
beide alleine waren. Gerade darum war er so wertvoll. Eines Tages konnte
Louis dem Drang, diese Zweisamkeit festzumachen, nicht widerstehen.
Er hatte sie beide aufgenommen. Sie hatten nur ein zwei
Biere getrunken, waren sich nahe gewesen und es klang bezaubernd.
Giorgia hätte niemals zugestimmt. Er gestand es ihr nicht. 
Die Erinnerung an die Aufnahme tröstete ihn. Auf seltsame Weise noch
immer. Obwohl er sie zusammen mit seinem Handy weggeworfen hatte.
Louis fielen die Worte ein.
«Liebe ist so dumm. Aber sie ist es wert, verstanden zu werden.» Er runzelte die Stirn und ein «pff» kam über seine Lippen.
\emph{Verstanden - } was eigentlich gab es zu verstehen? Sie hatten wohl
beide geglaubt, sich einig zu sein. Oder hatte er sich etwas vorgemacht? 
Mit den Monaten sangen sie den Song weniger, dann selten, und
schliesslich war er verschwunden. Zeitgleich mit der sich nähernden Unsicherheit.

\sectionbreak

Der sonnenhelle Tag schmerzte im scharfen Kontrast zu Louis’ Stimmung. Seine schweissnassen Kleider klebten.
Die Ausgelassenheit der Menschen um ihn herum war kaum zu ertragen. Die Trauer saugte ihm
verbissen seine Lebenskraft ab. Es war erniedrigend. Er fühlte sich aussen vor. Schlimmer noch, der Zustand war ihm bekannt.

Zum Teufel, er wusste nicht, wie Leben ging.

Vielleicht stand er endgültig vor dem Durchdrehen. Er dachte zu viel, zu kompliziert und zu versessen über Vergangenes nach. Louis hielt sich den Kopf, schloss die Augen und presste sein Gesicht in die Hände. Seine Gedanken erschienen ihm wie Häftlinge, die sich gegenseitig schikanierten, um die Sehnsucht nach Befreiung auszuhalten. 
Unvermittelt erschien Giorgia vor ihm. Diesmal wie eine Nahaufnahme auf Grossleinwand. Mit scharf abgegrenzten Umrissen.
Er sah ihre dunklen, offenen Augen, als würde sie gerade vor ihm stehen, gefährlich schwankend. Ihn ansehen und anschreien. Er brachte ihren sanften Blick schlicht nicht mit ihren wütend herausgepressten Worten zusammen. 

\sectionbreak

Wieder auf der Strasse kam ihm ein junges Paar entgegen. Sie lachten und riefen Louis ein «hi» zu. Er blieb stehen.

«Hi, habt ihr einen Moment?»

Die beiden drehten sich um. Sie sprachen englisch. Der Mann stellte seinen Gitarrenkoffer ab. Louis erklärte ihnen sein Anliegen.

«Wait,,» der Mann sprach zuerst, «geh ans Seeufer, oder noch
besser, zum «Helvetiaplatz», nicht weit von hier. Dort findest du
Künstler, auch Strassenmusiker. Vielleicht weiss einer was.»

«Seit letztem Jahr ist Strassenmusik auf dem Platz zugelassen und ein
Gewinn,» ergänzte die Frau, «grosser Platz, viel People und gute Einnahmen.»

Sie sahen vergnügt aus, erklärten ihm den Weg und zeigten Louis die
Richtung an.

«Easy, wirst du finden.»

Sie wünschten ihm Glück. 

Nach wenigen Minuten lag der Platz vor ihm. Ein betriebsamer, bunter Ort. 
Mit dem Verwaltungsgebäude, dem Volkshaus und dem Kanzleiareal umzäunten stattliche Gebäude den «Helvetiaplatz».
Louis sah sich neugierig um. Seine Hoffnung auf einen wertvollen Tipp wuchs.
Er durchschritt den Platz und blieb bei einem jungen Paar stehen.
Umringt von Erwachsenen und Kindern interpretierten sie Songs bekannter Musikgrössen. Er an der Gitarre. 
Sie spielte die Violine. Beide hatten ein Mic umgehängt und sangen gerade Bruce Springsteens \emph{«If I
Should Fall Behind»}. Sie hatten verblüffende Strahlkraft. Eine Schranke in Louis begann sich zu heben. Das mussten Profis sein.

\qrblock{Bruce Springsteen \emph{«If I Should Fall Behind»}}{media/QR-Codes/020_If_I_Should_Fall_Behind_-_Live_at_the_Point_Theatre,_Dublin,_Ireland_-_November_2006_Bruce_Springsteen_Spotify.png}
